\scanpage{3}
\diagram{\begin{tikzcd}[ampersand replacement=\&,column sep=large]
\text{Cat des systèmes de racines}\arrow[r,"\Phi"] \& \text{cat. des systèmes}
\end{tikzcd}}
\[
(W,R,\deleted{\unsure{\varphi}}I,\varphi:R\times I\to W)
\]
où $W$ est un groupe, $R$ un $W$-torseur à g., \inserted{$I$ ens. fini,} $\varphi:R\times I$ une \deleted{\unsure{de|$I$ $W$}} application

avec

a) $\varphi(wr,i)=\operatorname{int}(w)\varphi(r,\unsure{i'|i})$

b) $\forall r\in R,\quad i\mapsto\varphi(r,i)$ définit un système de générateurs de $W$ qui en fait \deleted{\unsure{de}} un \textquotedblleft groupe de Weyl épinglé\textquotedblright\ i.e.

c)
\[
\left\{\begin{aligned}\varphi(r,i)^2&=1\\(\varphi(r,i)\varphi(r,j))^{n_{ij}}&=1\quad(i\ne j)\end{aligned}\right.
\]
est une \uline{présentation} de $W$ (où $n_{ij}$ est l'ordre de $\varphi(r,i)\varphi(r,j)$), et

d) $I$ muni \deleted{\unsure{des}} \inserted{du} syst. des $n_{ij}$ ($n_{ij}=n_{ji}$) est un diagramme de Dynkin ($i$ lié à $j$ sss $n_{ij}\geq 3$ i.e. $n_{ij}\ne2$ i.e. $\varphi(r,i)$ ne commute pas à $\varphi(r,j)$)

\uline{NB} On a
\[
 R\xrightarrow{\psi}\ \text{Syst. gén admissibles de }W
 \quad\text{paramétrés par le diagramme de Dynkin }I
\]
Mais cette application, \deleted{\unsure{il se peut}} est \unsure{bien} \unsure{a)} compatible aux actions de $W$, n'\deleted{\unsure{est}}est pas \inserted{\unsure{tjrs|en gén.}} injective, \inserted{si $r\in R$, elle est inj.} \deleted{(si $I\ne\varnothing$), \unsure{car en fait} si $r,r'$ sont deux \textquotedblleft\unsure{repères}\textquotedblright\ correspondant à des syst. simples de racines opposés l'élément \inserted{$w$} du groupe de Weyl qui \unsure{amè|envoie}} ssi le stabilisateur de \unsure{$\psi$|$\psi(r)$|repères} ds $W$ est réduit à l'élément neutre. Or ce stabilisateur est évidemment le \uline{centre} de $W$ -- donc $\psi$ injective (permet d'interpréter $R$ comme un ens de syst. de gén. adm. de $W$ paramétrés par le diagramme de Dynkin $I$) \deleted{\illegible{1}} sss $Z(\unsure{W|w})=\{1\}$.

Le foncteur précédent est-il une équivalence de catégories ? Il est \inserted{\unsure{évid.\textsuperscript{t}|évid.|évid. fid.}} \unsure{ess\textsuperscript{t}|très|ess.} surjectif, est-il pl. fidèle ? Deux syst. de racines $\Sigma,\Sigma'$ tels que $\Phi(\Sigma)\sim\Phi(\Sigma')$ ont \deleted{\unsure{di}} gr. de Dynkin isom,
% archive sheet number: 2

